% 第10回　提案のブラッシュアップ（記入例・模範解答）
\session{10}{2}{12月2日（水）}{提案のブラッシュアップ}{AIによる批判的検証}{yes}

\goals{
  \item 企業が提案を見る観点（課題への適合、実現性、費用対効果、ブランド整合、運用負荷）を挙げられる
  \item AIに批判者の役割を与え、提案の弱点を体系的に洗い出せる
  \item 根拠となるデータの必要性を見極め、AIの数値ではなく一次資料で補強できる
}

\work[100mm]{準備}{提案の要点を1枚にまとめる}
\begin{wstab}{colspec={Q[l,wd=30mm,m] X[l,m]},row{1-Z}={ht=13mm}}
  \hc{企業の課題} & \af{20〜30代と子育て世帯の来店が少ない。オンラインで下調べを済ませた客に、店で何を提供するか} \\
  \hc{顧客（ペルソナ）} & \af{A：子連れで自分のペースで比べたい32歳／B：営業トークが苦手な初めての購入者24歳} \\
  \hc{コンセプト} & \af{「見るだけ、大歓迎。」子どもも自分も、自分のペースで過ごせる“公園ラウンジ”のディーラー} \\
  \hc{空間} & \af{展示車で囲んだ中央のキッズスペース、芝生風の床、半個室の商談ブース、ガラス越しの整備工場} \\
  \hc{体験（動線）} & \af{サインカードで「見ているだけ」を選べる。相談したいときだけカードで呼ぶ。待ち時間は子どもと過ごせる} \\
  \hc{期待される効果} & \af{若い世代・子育て世帯の来店のハードルを下げ、「相談したい」と思った客を商談につなげる} \\
  \hc{実現の道筋} & \af{①サインカードと案内サイン（すぐ試せる）　②家具の配置換え（改装中も可）　③床と照明の改装（段階的に）} \\
\end{wstab}

\work[80mm]{ワーク1・2}{企業側の視点で批判させ、対応を決める}
\begin{promptbox}[企業側の視点で批判させる]
あなたは自動車ディーラーを運営する会社の経営企画担当者です。次の店舗設計提案について、課題への適合、初期費用と運営費、実現の難しさ、ブランドイメージとの整合、スタッフの運用負荷、安全と法規の観点から、厳しく問題点を指摘してください。指摘は重要な順に並べてください。提案：（要約）
\end{promptbox}
\instr{分類：Ａ＝対応する／Ｂ＝根拠を示して反論する／Ｃ＝割り切る（理由を書く）}
\begin{wstabh}{colspec={X[3,l,m] Q[l,wd=18mm,m] Q[c,wd=10mm,m] X[3,l,m]},row{2-Z}={ht=12mm}}
  AIの指摘（重要な順） & 観点 & 分類 & 具体的な対応・反論の根拠・割り切る理由 \\
  \af{「話しかけない」接客では、商談につながらず売上が下がるのでは} & \af{課題への適合} & \ans{Ｂ} & \af{課題は「来店のハードル」。セルフ方式を取り入れた他業種の事例を示し、相談の質が上がることを根拠にする} \\
  \af{展示ゾーン中央のキッズスペースは、車の出し入れの際に危険} & \af{安全と法規} & \ans{Ａ} & \af{展示車は固定し、出し入れは閉店後に限る。キッズスペースに柵を設ける} \\
  \af{見守り担当の人件費がかかる} & \af{運用負荷} & \ans{Ａ} & \af{休日の混雑時だけ配置。平日は受付と兼務する} \\
  \af{芝生風の床はメーカーのブランド基準に合わないかもしれない} & \af{ブランド整合} & \ans{Ａ} & \af{企業担当者に基準を確認する。合わない場合はキッズスペースだけに限る} \\
  \af{改装費用に見合う来店増の根拠がない} & \af{費用対効果} & \ans{Ｃ} & \af{来店数の予測は根拠となるデータがないため示さない。費用の小さい施策から段階的に試す提案にする} \\
  \af{サインカードの効果が不明} & \af{実現の難しさ} & \ans{Ｃ} & \af{まず1か月試し、来店客の反応を記録する形で提案する} \\
\end{wstabh}

\work[60mm]{根拠}{根拠が示されていない主張を見つける}
\begin{promptbox}[根拠の不足を見つける]
この提案の中で、根拠（データ、事例、調査）が示されていない主張を列挙し、それぞれについて、どんな一次資料で裏づけられそうかを提案してください。数値やデータそのものは生成しないでください。
\end{promptbox}
\begin{wstabh}{colspec={X[3,l,m] X[3,l,m] Q[l,wd=18mm,m]},row{2-Z}={ht=11mm}}
  根拠が不足している主張 & 裏づけに使う一次資料（見つけたら出典） & 担当 \\
  \af{若い世代は販売員との対話を負担に感じている} & \af{前半で自分が行ったアンケート、公的調査や業界団体の調査（出典と頁を記入）} & \af{自分} \\
  \af{子育て世帯は子どもを見ながら車を選びたい} & \af{子育て中の家族への聞き取り（5組）、企業担当者の話} & \af{メンバーＢ} \\
  \af{セルフ方式でも相談の質は下がらない} & \af{セルフ方式を取り入れた他業種の事例（企業の公開資料・報道）} & \af{メンバーＣ} \\
  \af{段階的な改装なら営業を続けられる} & \af{企業担当者への確認（第9回のリスト）} & \af{メンバーＤ} \\
\end{wstabh}

\work[60mm]{中間発表}{各グループ5分。受けた指摘と、他グループへのコメント}
\begin{wstabh}{colspec={X[l,m] X[l,m]},row{2-Z}={ht=12mm}}
  自分たちが受けた質問・指摘 & 対応方針 \\
  \af{整備で来る客（既存客）にとって、店が変わることのデメリットは？} & \af{整備の待合を静かなエリアとして残す。ペルソナに既存客を1人追加して検証する} \\
  \af{サインカードを持たない客（ふらっと来た客）はどうするのか} & \af{入口のサインで説明し、受付で一言だけ案内する} \\
  \af{「公園」というコンセプトが外観から伝わるか} & \af{外から見えるガラス面にキッズスペースを置く案を検討する} \\
\end{wstabh}
\abox[他グループへのコメント（グループ名と内容）]{28mm}{グループA：夜だけ開く店は面白いが、スタッフの勤務時間の負担について説明があるとよい。\newline
グループD：整備待ちを体験に変える案は、既存客を大切にしていて説得力がある。体験の内容をもっと具体的に聞きたい。}

\begin{nexttime}
  \item 提案の改訂版を作る（指摘への対応を明記）
  \item 発表の構成案のたたき台を作る
\end{nexttime}
\abox[改訂のポイント（何をどう変えたか）]{30mm}{① 安全：展示車の固定とキッズスペースの柵を追加。② 運用：見守り担当は休日のみ。③ 根拠：来店の負担感はアンケートと公的調査で示し、来店数の予測は示さない。④ 既存客：整備の静かな待合を残す。⑤ 実現の道筋：費用の小さい施策から3段階で試す形に整理。}
\aimemoA{企業側の立場からの批判と、根拠が不足している主張の洗い出し}
{「経営企画担当者として、6つの観点から厳しく問題点を、重要な順に」「根拠が示されていない主張と、裏づけに使えそうな一次資料を。数値は生成しないで」}
{6つの指摘を対応・反論・割り切るに分類した。AIが例として挙げた来店増の数値は、出典がないため使わなかった。}
{裏づけは自分たちのアンケート、聞き取り、企業の公開資料で集めた。中間発表で他グループの指摘も受けた}

\begin{point}
  \item 指摘を「対応・反論・割り切る」に分類できているかを見る。すべてに対応しようとすると提案がぶれる。割り切る場合は理由が書けていればよい。
  \item AIに数値を出させて根拠にしていないかを必ず確認する。「来店が○\%増える」のような数値は、出典がなければ使わせない。
  \item 中間発表の指摘への対応方針が、改訂のポイントに反映されているかを見る。
  \item 他グループへのコメントは、良い点と質問をセットで書けているかを評価する。
\end{point}
